% Unit 7 review sheet -- 2 pages.
\documentclass{apchem-worksheet}

\apcunit{7}
\apcunittitle{Equilibrium}
\apcdoctitle{Review Sheet}
\apcsubtitle{Everything on the exam traces to this page \textbullet\ exam Mon Jan 25}

\begin{document}
\apctitleblock

\section*{\sffamily The five tools}

\renewcommand{\arraystretch}{1.45}
\noindent\begin{tabular}{@{}p{3.9cm}p{9.8cm}@{}}
\toprule
\textbf{$Q$ versus $K$} & $Q<K$ shift right, $Q=K$ at equilibrium, $Q>K$
  shift left. Works for \emph{any} mixture, equilibrium or not --- the one
  tool that answers ``which way?''\\
\textbf{Writing $K$} & products over reactants, each raised to its
  coefficient. Omit pure solids and pure liquids. \ce{H2O(l)} out,
  \ce{H2O(g)} in\\
\textbf{ICE} & change row follows the coefficients; approximate only if
  justified, and \emph{report the percentage}. Perfect squares beat the
  quadratic formula\\
\textbf{Manipulating $K$} & reverse $\to 1/K$; $\times n$ coefficients
  $\to K^n$; add reactions $\to K_1K_2$. Constants multiply where
  enthalpies add\\
\textbf{Le Ch\^atelier} & the system opposes the change. Only
  \textbf{temperature} alters the value of $K$\\
\bottomrule
\end{tabular}

\section*{\sffamily The disturbance table}

\begin{center}
\renewcommand{\arraystretch}{1.15}
\begin{tabular}{@{}lll@{}}
\toprule
\textbf{Disturbance} & \textbf{Shift} & \textbf{$K$?} \\
\midrule
add reactant / remove product & right & unchanged \\
add product / remove reactant & left & unchanged \\
compress (smaller volume) & toward fewer moles of gas & unchanged \\
add inert gas at constant $V$ & none & unchanged \\
add catalyst & none & unchanged \\
raise $T$, forward endothermic & right & \textbf{increases} \\
raise $T$, forward exothermic & left & \textbf{decreases} \\
\bottomrule
\end{tabular}
\end{center}

\section*{\sffamily Solubility, in four lines}

\begin{itemize}[leftmargin=*,itemsep=0.18em]
  \item \ce{MX}: $K_{sp} = s^2$. \ce{MX2} or \ce{M2X}:
        $K_{sp} = 4s^3$.
  \item $K_{sp}$ is a \emph{constant}; solubility is a \emph{position}
        that changes with what else is dissolved.
  \item Common ion $\Rightarrow$ $Q$ exceeds $K_{sp}$ $\Rightarrow$ less
        dissolves. $K_{sp}$ itself never moves.
  \item Compare $K_{sp}$ values directly \textbf{only} when the salts
        release the same number of ions.
\end{itemize}

\section*{\sffamily The traps, one line each}

\begin{itemize}[leftmargin=*,itemsep=0.22em]
  \item Saying equilibrium means equal concentrations (it means equal
        rates).
  \item Including a pure solid or liquid in $K$.
  \item Forgetting an exponent --- $(0.30)^3$ is $0.027$, not $0.90$.
  \item Using the small-$x$ approximation without checking, or checking and
        ignoring the result.
  \item Claiming a catalyst or an inert gas shifts an equilibrium.
  \item Saying a concentration change alters $K$.
  \item Comparing $K_{sp}$ for salts with different ion counts.
\end{itemize}

\section*{\sffamily Self-check --- 12 questions, exam conditions}

\begin{enumerate}[leftmargin=*,itemsep=0.38em]
  \item $K$ for $\ce{CaCO3(s) <=> CaO(s) + CO2(g)}$:
        \answerblank[2.6cm]{$[\ce{CO2}]$}
  \item $[\ce{H2}]=0.20$, $[\ce{I2}]=0.40$, $[\ce{HI}]=0.60$. Find $K$:
        \answerblank[2cm]{4.5}
  \item $K = 4.0$; give $K$ for the reverse:
        \answerblank[2cm]{0.25}
  \item $K = 4.0$; give $K$ with coefficients doubled:
        \answerblank[2cm]{16}
  \item $K = 4.5$ and $Q = 9.0$. Direction?
        \answerblank[2.4cm]{shifts left}
  \item Threshold for the small-$x$ approximation:
        \answerblank[1.6cm]{5\%}
  \item Adding argon at constant volume:
        \answerblank[2.4cm]{no shift}
  \item Forward reaction exothermic; raise $T$. What happens to $K$?
        \answerblank[2.4cm]{decreases}
  \item Molar solubility of \ce{AgCl}, $K_{sp} = 1.6\times10^{-10}$:
        \answerblank[3cm]{$1.3\times10^{-5}$~M}
  \item Solubility of \ce{AgCl} in \SI{0.10}{\Molar} \ce{NaCl}:
        \answerblank[3cm]{$1.6\times10^{-9}$~M}
  \item $K_{sp}$ for an \ce{MX2} salt in terms of $s$:
        \answerblank[1.8cm]{$4s^3$}
  \item Only what change alters $K$?
        \answerblank[2.6cm]{temperature}
\end{enumerate}

\begin{apcnote}[How to study for this exam]
This is a 44-point exam, one of the two largest in the course, so budget
accordingly. Redo WS 5's FRQ cold --- the long free-response is its twin in
shape, and it deliberately runs the full chain: expression $\to$ ICE $\to$
solve $\to$ verify $\to$ $Q$ versus $K$ $\to$ effect of a disturbance.

Two habits are worth more than any extra content. First, whenever you write
``shifts right,'' add \emph{because $Q < K$} --- the quantitative
justification is usually the point that the qualitative one is not. Second,
whenever you approximate, write the percentage down. Both are cheap, and
both are routinely omitted.
\end{apcnote}

\end{document}
